% GENERATED FILE. Edit canonical lessons, then run npm run generate:books.
% canonical-source-hash: fnv1a64:d82c53e80d9ca32a
% canonical-lessons: HI-C508-intarvyu, HI-C508-niyukti, HI-C508-avedan, HI-C508-dakhila, HI-C508-vetan

\chapter{Looking for a Job}
\label{ch:hi-looking-for-a-job}
\hlchaptermodality{508}

\begin{chapteropening}
\textbf{By the end of this chapter:} \emph{I can apply for a job or a place and say what I can do.}

Complete the last lesson of chapter 508: I can apply for a job or a place and say what I can do.
\end{chapteropening}

\section[\texorpdfstring{\hi{इंटरव्यू} (iṇṭarvyū)}{इंटरव्यू (iṇṭarvyū)}]{\hi{इंटरव्यू} (iṇṭarvyū) — an interview (for a job or a place)}

\label{lesson:HI-C508-intarvyu}



\begin{quote}
Before the new one: say the Hindi for skill, then the Hindi for talent.
\end{quote}

\subsection*{You'll want to know: \hi{इंटरव्यू}}
\textbf{\hi{इंटरव्यू}} — \emph{iṇṭarvyū} — \textquotedblleft{}an interview (for a job or a place)\textquotedblright{}.

\begin{quote}
\hi{कल} \hi{सुबह} \hi{दस} \hi{बजे} \hi{मेरा} \hi{इंटरव्यू} \hi{है।}
\end{quote}

\begin{quote}
\emph{Kal subah das baje merā iṇṭarvyū hai.}
\end{quote}

\begin{quote}
My interview is at ten tomorrow morning.
\end{quote}

Masculine.

\begin{cousinweb}
\textbf{\hi{इंटरव्यू}} (\emph{iṇṭarvyū}) is a loan from English \emph{interview}.
\end{cousinweb}

\begin{grammarlens}[title={What you've built}]
One more word for describing what you lived and felt.
\end{grammarlens}

\subsection*{Your turn}
\begin{itemize}
\raggedright
  \item \emph{Say it:} \emph{iṇṭarvyū}
  \item \emph{Say it:} \emph{iṇṭarvyū}, once more
  \item \emph{Say it:} \emph{Kal subah das baje merā iṇṭarvyū hai.} — My interview is at ten tomorrow morning.
  \item \emph{Recall:} say the Hindi for skill, then the Hindi for talent, then say \emph{iṇṭarvyū} again
\end{itemize}

\begin{tcolorbox}[breakable,colback=teal!4,colframe=teal!35!black,title={Before you move on}]
What does \hi{इंटरव्यू} mean? (\textquotedblleft{}An interview (for a job or a place)\textquotedblright{}.) Say it once more.
\end{tcolorbox}
\section[\texorpdfstring{\hi{नियुक्ति} (niyukti)}{नियुक्ति (niyukti)}]{\hi{नियुक्ति} (niyukti) — an appointment (to a post)}

\label{lesson:HI-C508-niyukti}



\begin{quote}
Before the new one: say the Hindi for talent, then the Hindi for an interview (for a job or a place).
\end{quote}

\subsection*{You'll want to know: \hi{नियुक्ति}}
\textbf{\hi{नियुक्ति}} — \emph{niyukti} — \textquotedblleft{}an appointment (to a post)\textquotedblright{}.

\begin{quote}
\hi{अगले} \hi{साल} \hi{मेरी} \hi{नियुक्ति} \hi{एक} \hi{नए} \hi{स्कूल} \hi{में} \hi{होगी।}
\end{quote}

\begin{quote}
\emph{Agle sāl merī niyukti ek nae skūl meṁ hogī.}
\end{quote}

\begin{quote}
Next year I will be appointed to a new school.
\end{quote}

Feminine.

\begin{cousinweb}
\textbf{\hi{नियुक्ति}} (\emph{niyukti}) is a tatsama, taken whole from Sanskrit \emph{niyukti}, \textquotedblleft{}a charge, an office\textquotedblright{} (Monier-Williams).
\end{cousinweb}

\begin{grammarlens}[title={What you've built}]
One more word for describing what you lived and felt.
\end{grammarlens}

\subsection*{Your turn}
\begin{itemize}
\raggedright
  \item \emph{Say it:} \emph{niyukti}
  \item \emph{Say it:} \emph{niyukti}, once more
  \item \emph{Say it:} \emph{Agle sāl merī niyukti ek nae skūl meṁ hogī.} — Next year I will be appointed to a new school.
  \item \emph{Recall:} say the Hindi for talent, then the Hindi for an interview (for a job or a place), then say \emph{niyukti} again
\end{itemize}

\begin{tcolorbox}[breakable,colback=teal!4,colframe=teal!35!black,title={Before you move on}]
What does \hi{नियुक्ति} mean? (\textquotedblleft{}An appointment (to a post)\textquotedblright{}.) Say it once more.
\end{tcolorbox}
\section[\texorpdfstring{\hi{आवेदन} (āvedan)}{आवेदन (āvedan)}]{\hi{आवेदन} (āvedan) — an application (formal)}

\label{lesson:HI-C508-avedan}



\begin{quote}
Before the new one: say the Hindi for an interview (for a job or a place), then the Hindi for an appointment (to a post).
\end{quote}

\subsection*{You'll want to know: \hi{आवेदन}}
\textbf{\hi{आवेदन}} — \emph{āvedan} — \textquotedblleft{}an application (formal)\textquotedblright{}.

\begin{quote}
\hi{मैंने} \hi{नौकरी} \hi{के} \hi{लिए} \hi{आवेदन} \hi{भेजा।}
\end{quote}

\begin{quote}
\emph{Maiṁne naukrī ke liye āvedan bhejā.}
\end{quote}

\begin{quote}
I sent an application for the job.
\end{quote}

Masculine.

\begin{cousinweb}
\textbf{\hi{आवेदन}} (\emph{āvedan}) is a tatsama, taken whole from Sanskrit \emph{āvedana}, \textquotedblleft{}announcing, informing\textquotedblright{} (Monier-Williams). Hindi uses it for a written request, an application.
\end{cousinweb}

\begin{grammarlens}[title={What you've built}]
One more word for describing what you lived and felt.
\end{grammarlens}

\subsection*{Your turn}
\begin{itemize}
\raggedright
  \item \emph{Say it:} \emph{āvedan}
  \item \emph{Say it:} \emph{āvedan}, once more
  \item \emph{Say it:} \emph{Maiṁne naukrī ke liye āvedan bhejā.} — I sent an application for the job.
  \item \emph{Recall:} say the Hindi for an interview (for a job or a place), then the Hindi for an appointment (to a post), then say \emph{āvedan} again
\end{itemize}

\begin{tcolorbox}[breakable,colback=teal!4,colframe=teal!35!black,title={Before you move on}]
What does \hi{आवेदन} mean? (\textquotedblleft{}An application (formal)\textquotedblright{}.) Say it once more.
\end{tcolorbox}
\section[\texorpdfstring{\hi{दाख़िला} (dākhilā)}{दाख़िला (dākhilā)}]{\hi{दाख़िला} (dākhilā) — admission (to a school or course)}

\label{lesson:HI-C508-dakhila}



\begin{quote}
Before the new one: say the Hindi for an appointment (to a post), then the Hindi for an application (formal).
\end{quote}

\subsection*{You'll want to know: \hi{दाख़िला}}
\textbf{\hi{दाख़िला}} — \emph{dākhilā} — \textquotedblleft{}admission (to a school or course)\textquotedblright{}.

\begin{quote}
\hi{मेरे} \hi{भाई} \hi{को} \hi{बड़े} \hi{स्कूल} \hi{में} \hi{दाख़िला} \hi{मिला।}
\end{quote}

\begin{quote}
\emph{Mere bhāī ko baṛe skūl meṁ dākhilā milā.}
\end{quote}

\begin{quote}
My brother got a place at the big school.
\end{quote}

Masculine.

\begin{cousinweb}
\textbf{\hi{दाख़िला}} (\emph{dākhilā}) is usually traced to Arabic \emph{dākhila}, \textquotedblleft{}an entry\textquotedblright{}, through Persian; Hindi uses it for being let into a school or a course.
\end{cousinweb}

\begin{grammarlens}[title={What you've built}]
One more word for describing what you lived and felt.
\end{grammarlens}

\subsection*{Your turn}
\begin{itemize}
\raggedright
  \item \emph{Say it:} \emph{dākhilā}
  \item \emph{Say it:} \emph{dākhilā}, once more
  \item \emph{Say it:} \emph{Mere bhāī ko baṛe skūl meṁ dākhilā milā.} — My brother got a place at the big school.
  \item \emph{Recall:} say the Hindi for an appointment (to a post), then the Hindi for an application (formal), then say \emph{dākhilā} again
\end{itemize}

\begin{tcolorbox}[breakable,colback=teal!4,colframe=teal!35!black,title={Before you move on}]
What does \hi{दाख़िला} mean? (\textquotedblleft{}Admission (to a school or course)\textquotedblright{}.) Say it once more.
\end{tcolorbox}
\section[\texorpdfstring{\hi{वेतन} (vetan)}{वेतन (vetan)}]{\hi{वेतन} (vetan) — a salary, pay (formal)}

\label{lesson:HI-C508-vetan}



\begin{quote}
Before the new one: say the Hindi for an application (formal), then the Hindi for admission (to a school or course).
\end{quote}

\subsection*{You'll want to know: \hi{वेतन}}
\textbf{\hi{वेतन}} — \emph{vetan} — \textquotedblleft{}a salary, pay (formal)\textquotedblright{}.

\begin{quote}
\hi{इस} \hi{नौकरी} \hi{में} \hi{वेतन} \hi{अच्छा} \hi{है}, \hi{लेकिन} \hi{काम} \hi{बहुत} \hi{है।}
\end{quote}

\begin{quote}
\emph{Is naukrī meṁ vetan acchā hai, lekin kām bahut hai.}
\end{quote}

\begin{quote}
This job pays well, but there is a lot of work.
\end{quote}

Masculine.

A formal word; everyday speech also says \hi{तनख़्वाह}.

\begin{cousinweb}
\textbf{\hi{वेतन}} (\emph{vetan}) is a tatsama, taken whole from Sanskrit \emph{vetana}, \textquotedblleft{}wages, hire, salary\textquotedblright{} (Monier-Williams).
\end{cousinweb}

\begin{grammarlens}[title={What you've built}]
One more word for describing what you lived and felt.
\end{grammarlens}

\subsection*{Your turn}
\begin{itemize}
\raggedright
  \item \emph{Say it:} \emph{vetan}
  \item \emph{Say it:} \emph{vetan}, once more
  \item \emph{Say it:} \emph{Is naukrī meṁ vetan acchā hai, lekin kām bahut hai.} — This job pays well, but there is a lot of work.
  \item \emph{Recall:} say the Hindi for an application (formal), then the Hindi for admission (to a school or course), then say \emph{vetan} again
\end{itemize}

\begin{tcolorbox}[breakable,colback=teal!4,colframe=teal!35!black,title={Before you move on}]
What does \hi{वेतन} mean? (\textquotedblleft{}A salary, pay (formal)\textquotedblright{}.) Say it once more.
\end{tcolorbox}
